\documentclass[11pt]{article}

% -------------------- Packages --------------------
\usepackage[a4paper,margin=1in]{geometry}
\usepackage{hyperref}
\usepackage{enumitem}
\usepackage{titlesec}
\usepackage{fancyhdr}
\usepackage{xcolor}
\usepackage{listings}
% listings has no built-in yaml or Docker language; define them (plain, no highlighting) so latexmk builds without errors.
\lstdefinelanguage{yaml}{}
\lstdefinelanguage{Docker}{}
\usepackage{courier}

% -------------------- Hyperref Setup --------------------
\hypersetup{
  colorlinks=true,
  linkcolor=blue,
  urlcolor=blue,
  citecolor=blue
}

% -------------------- Header/Footer --------------------
\pagestyle{fancy}
\fancyhf{}
\lhead{MCV Lab}
\rhead{Cloud ML Deployment Infrastructure}
\cfoot{\thepage}

% -------------------- Section Formatting --------------------
\titleformat{\section}{\Large\bfseries}{\thesection}{0.8em}{}
\titleformat{\subsection}{\normalsize\bfseries}{\thesubsection}{0.8em}{}
\titleformat{\subsubsection}{\normalsize\bfseries}{\thesubsubsection}{0.8em}{}

\titlespacing*{\section}{0pt}{1.5em}{0.8em}
\titlespacing*{\subsection}{0pt}{1.0em}{0.5em}
\titlespacing*{\subsubsection}{0pt}{0.8em}{0.4em}

% -------------------- Global List Formatting --------------------
\setlist[itemize]{
  label={-},
  leftmargin=2em,
  labelsep=0.6em,
  itemsep=0.2em
}
\setlist[itemize,2]{itemsep=0.15em}
\setlist[enumerate]{leftmargin=2em,itemsep=0.15em}

% -------------------- Listings / Code Style --------------------
\definecolor{codegray}{rgb}{0.95,0.95,0.95}
\definecolor{commentgreen}{rgb}{0.0,0.5,0.0}

\lstset{
  basicstyle=\footnotesize\ttfamily,
  backgroundcolor=\color{codegray},
  frame=single,
  breaklines=true,
  showstringspaces=false,
  keywordstyle=\color{blue},
  commentstyle=\color{commentgreen},
  stringstyle=\color{red},
  columns=fullflexible,
  keepspaces=true,
  tabsize=2,
  xleftmargin=0.02\textwidth,
  xrightmargin=0.02\textwidth
}

% -------------------- Document Metadata --------------------
\title{\textbf{Medical Computer Vision Lab \\ \small{
Cloud Platforms (AWS, Azure, GCP) and AI/ML Model Deployment Infrastructure: Lab Standard}}}
\author{Maintainer: Ashery Mbilinyi}
\date{Version 1.0 : \today}

\begin{document}
\maketitle

\tableofcontents
\newpage

\section{Purpose}
This document explains cloud platforms (AWS, Azure, GCP) in the context of AI/ML workloads, and defines practical lab standards for \textbf{model deployment infrastructure} in the cloud.
The goal is to ensure our cloud deployments are:
\begin{itemize}
  \item \textbf{Reproducible}: consistent environments and versioned model artifacts
  \item \textbf{Secure}: least-privilege access and safe secret handling
  \item \textbf{Scalable}: supports increased traffic, batching, and multi-replica serving
  \item \textbf{Observable}: logs, metrics, and basic monitoring are always present
  \item \textbf{Rollback-ready}: model versions can be promoted and reverted safely
\end{itemize}

\section{What ``Cloud Platforms'' Provide for ML}
Cloud platforms provide standardized building blocks:

\subsection{Compute}
Where code runs:
\begin{itemize}
  \item CPU VMs for web services and light inference
  \item GPU VMs for training and GPU inference
  \item managed container compute (serverless containers)
\end{itemize}

\subsection{Storage}
Where data and artifacts live:
\begin{itemize}
  \item datasets, checkpoints, logs
  \item model artifacts (e.g., \texttt{.pt}, \texttt{.onnx}, tokenizers)
\end{itemize}

\subsection{Networking}
How services are exposed and connected:
\begin{itemize}
  \item load balancers, DNS, HTTPS termination
  \item private networks (VPC/VNet), firewall rules
\end{itemize}

\subsection{Identity \& Access Management (IAM)}
Who can access resources:
\begin{itemize}
  \item service accounts/roles
  \item permission policies
  \item principle of least privilege
\end{itemize}

\section{What ``Model Deployment Infrastructure'' Means}
Model deployment infrastructure is the full system required to serve a model reliably:
\begin{itemize}
  \item \textbf{Model packaging}: containerize inference + preprocessing code
  \item \textbf{Artifact storage}: store model files and metadata
  \item \textbf{Serving layer}: API endpoint (REST/gRPC)
  \item \textbf{Scaling}: autoscaling + load balancing + batching strategies
  \item \textbf{Monitoring}: latency, error rates, drift signals (when available)
  \item \textbf{Versioning}: model registry + rollout/rollback
\end{itemize}

\section{Equivalent Services: AWS vs Azure vs GCP}

\subsection{Core Infrastructure Mapping (Quick Reference)}
\begin{itemize}
  \item \textbf{Compute VMs}: AWS EC2 / Azure VMs / GCP Compute Engine
  \item \textbf{Storage}: AWS S3 / Azure Blob Storage / GCP Cloud Storage
  \item \textbf{Container Registry}: AWS ECR / Azure ACR / GCP Artifact Registry
  \item \textbf{Kubernetes}: AWS EKS / Azure AKS / GCP GKE
  \item \textbf{Managed ML Platform}: AWS SageMaker / Azure ML / GCP Vertex AI
  \item \textbf{Secrets}: AWS Secrets Manager / Azure Key Vault / GCP Secret Manager
\end{itemize}

\section{Three Deployment Patterns (Most Common)}

\subsection{Pattern A: Managed ML Endpoints (Fastest to Production)}
Use cloud-managed model serving when you want simplicity and minimal DevOps.

Examples:
\begin{itemize}
  \item AWS SageMaker Endpoints
  \item Azure ML Online Endpoints
  \item GCP Vertex AI Endpoints
\end{itemize}

\textbf{Pros:}
\begin{itemize}
  \item fastest setup
  \item built-in scaling/monitoring integrations
  \item standardized deployment workflow
\end{itemize}

\textbf{Cons:}
\begin{itemize}
  \item less control over runtime details
  \item cost can be higher for always-on endpoints
\end{itemize}

\subsection{Pattern B: Kubernetes + Containers (Most Flexible)}
Deploy a containerized inference API (FastAPI) on Kubernetes.

\textbf{Pros:}
\begin{itemize}
  \item maximum control and portability
  \item strong scaling patterns and reliability
  \item standard infra across clouds
\end{itemize}

\textbf{Cons:}
\begin{itemize}
  \item operational complexity
  \item requires cluster management knowledge
\end{itemize}

\subsection{Pattern C: Serverless Containers (Great Middle Ground)}
Run inference as a container on managed serverless container compute.

Examples:
\begin{itemize}
  \item AWS: ECS Fargate (serverless containers)
  \item Azure: Container Apps
  \item GCP: Cloud Run
\end{itemize}

\textbf{Pros:}
\begin{itemize}
  \item minimal ops
  \item easy autoscaling
  \item good for CPU inference and many API workloads
\end{itemize}

\textbf{Cons:}
\begin{itemize}
  \item GPU inference support varies by platform/region and can be costly
  \item cold starts may matter for low-traffic services
\end{itemize}

\section{Lab Standard: Recommended Default Choices}
\subsection{Default Recommendations}
\begin{itemize}
  \item \textbf{Internal lab demos}: Docker + FastAPI on a single VM (CPU/GPU as needed)
  \item \textbf{Simple production-style API}: serverless containers (Cloud Run / Container Apps / Fargate)
  \item \textbf{Research platform at scale}: Kubernetes (GKE/EKS/AKS) or Managed Endpoints
\end{itemize}

\subsection{GPU Inference Guidance}
Scaling GPU inference requires careful control:
\begin{itemize}
  \item avoid too many workers competing for the same GPU memory
  \item prefer batching where latency allows
  \item consider queue-based designs for heavy workloads
\end{itemize}

\section{Deployment Building Blocks (What Every Cloud Deployment Needs)}
\subsection{(1) Model Artifacts and Metadata}
We store:
\begin{itemize}
  \item model weights (\texttt{model.pt} / \texttt{model.onnx})
  \item preprocessing config (normalization, transforms)
  \item label map / class definitions
  \item model card metadata (dataset, metrics, caveats)
\end{itemize}

\subsection{(2) Containerized Inference Service}
We serve inference via a container image:
\begin{itemize}
  \item FastAPI app exposes \texttt{/health}, \texttt{/ready}, \texttt{/version}, \texttt{/predict}
  \item container pulls model artifact at startup or includes it in the image (decision depends on update frequency)
\end{itemize}

\subsection{(3) Traffic Management}
\begin{itemize}
  \item HTTPS/TLS termination (managed load balancer or ingress)
  \item routing and versioning (v1/v2)
  \item optional canary rollouts
\end{itemize}

\subsection{(4) Observability}
Minimum:
\begin{itemize}
  \item structured logs
  \item latency and error rate metrics
\end{itemize}
Recommended:
\begin{itemize}
  \item tracing (distributed traces for microservices)
  \item drift monitoring when feasible
\end{itemize}

\subsection{(5) Security}
\begin{itemize}
  \item IAM roles for services (least privilege)
  \item secrets stored in managed secret store (not in code or Docker images)
  \item private networking when possible
\end{itemize}

\section{Coding Examples}

\subsection{Minimal FastAPI Inference App (Cloud-Friendly)}
\begin{lstlisting}[language=Python]
# app.py
from __future__ import annotations

import os
from typing import List, Optional

from fastapi import FastAPI, HTTPException
from pydantic import BaseModel, Field

APP_VERSION = os.getenv("APP_VERSION", "dev")
MODEL_VERSION = os.getenv("MODEL_VERSION", "unknown")

app = FastAPI(title="MCV Cloud Inference API", version=APP_VERSION)

MODEL_LOADED = False

class PredictRequest(BaseModel):
    x: List[float] = Field(..., min_length=1)
    return_probs: bool = False

class PredictResponse(BaseModel):
    y: int
    probs: Optional[List[float]] = None
    model_version: str

def load_model() -> None:
    # Replace with: download weights from S3/Blob/GCS or mount volume.
    global MODEL_LOADED
    MODEL_LOADED = True

def predict_impl(x: List[float]) -> tuple[int, List[float]]:
    score = sum(x)
    y = 1 if score > 0 else 0
    probs = [0.2, 0.8] if y == 1 else [0.8, 0.2]
    return y, probs

@app.on_event("startup")
def startup_event():
    load_model()

@app.get("/health")
def health():
    return {"status": "ok"}

@app.get("/ready")
def ready():
    if not MODEL_LOADED:
        raise HTTPException(status_code=503, detail="Model not loaded")
    return {"status": "ready"}

@app.get("/version")
def version():
    return {"app_version": APP_VERSION, "model_version": MODEL_VERSION}

@app.post("/predict", response_model=PredictResponse)
def predict(req: PredictRequest):
    if not MODEL_LOADED:
        raise HTTPException(status_code=503, detail="Model not loaded")

    y, probs = predict_impl(req.x)
    return PredictResponse(
        y=y,
        probs=probs if req.return_probs else None,
        model_version=MODEL_VERSION,
    )
\end{lstlisting}

\subsection{Dockerfile (Inference Service)}
\begin{lstlisting}[language=Docker]
FROM python:3.10-slim

WORKDIR /app

RUN pip install --no-cache-dir fastapi uvicorn[standard] pydantic

COPY app.py /app/app.py

EXPOSE 8080
CMD ["uvicorn", "app:app", "--host", "0.0.0.0", "--port", "8080"]
\end{lstlisting}

\subsection{Local Run}
\begin{lstlisting}[language=bash]
docker build -t mcv-inference:latest .
docker run -p 8080:8080 \
  -e APP_VERSION=dev \
  -e MODEL_VERSION=v0 \
  mcv-inference:latest
\end{lstlisting}

\subsection{Example: Deploy to GCP Cloud Run (Conceptual)}
\begin{lstlisting}[language=bash]
# Build and push to Artifact Registry (names simplified)
gcloud auth configure-docker
docker tag mcv-inference:latest REGION-docker.pkg.dev/PROJECT/REPO/mcv-inference:latest
docker push REGION-docker.pkg.dev/PROJECT/REPO/mcv-inference:latest

# Deploy
gcloud run deploy mcv-inference \
  --image REGION-docker.pkg.dev/PROJECT/REPO/mcv-inference:latest \
  --platform managed \
  --region REGION \
  --allow-unauthenticated
\end{lstlisting}

\subsection{Example: Deploy to AWS (Conceptual ECS/Fargate)}
\begin{lstlisting}[language=bash]
# Build -> push to ECR -> create ECS service (high-level sketch)
# 1) Push image to ECR
# 2) Define task definition (CPU/mem/env vars)
# 3) Create service behind load balancer
# 4) Enable autoscaling rules
\end{lstlisting}

\subsection{Example: Deploy to Azure (Conceptual Container Apps)}
\begin{lstlisting}[language=bash]
# Build -> push to ACR -> deploy to Azure Container Apps (high-level sketch)
# 1) Push image to ACR
# 2) Create Container App with ingress enabled
# 3) Configure autoscaling rules
\end{lstlisting}

\section{CI/CD for Cloud Deployments (Minimal Standard)}
A typical pipeline:
\begin{enumerate}
  \item CI: lint + unit tests + inference smoke test
  \item Build: Docker image
  \item Push: image to registry (ECR/ACR/Artifact Registry)
  \item Deploy: update service (Cloud Run / ECS / Container Apps / K8s)
\end{enumerate}

\subsection{Minimal GitHub Actions Sketch (Build Only)}
\begin{lstlisting}[language=yaml]
name: Build

on:
  push:
    branches: [ "main" ]

jobs:
  docker-build:
    runs-on: ubuntu-latest
    steps:
      - uses: actions/checkout@v4
      - name: Build
        run: docker build -t mcv-inference:latest .
\end{lstlisting}

\section{Checklist (Code Review)}
Before merging a cloud deployment project:
\begin{itemize}
  \item Inference API implements \texttt{/health}, \texttt{/ready}, \texttt{/version}, \texttt{/predict}
  \item Dockerfile exists and builds successfully
  \item Model version is exposed (\texttt{/version}) and traceable to an artifact
  \item Secrets are not in the repo or image
  \item Logging exists and is structured enough to debug failures
  \item Scaling approach is documented (replicas/autoscaling/batching)
\end{itemize}

\end{document}
